% !TeX program = xelatex
% File tu chua: khong can .bib, hinh anh hay file \input ben ngoai.
% Bien dich: xelatex BAO_CAO_KHOA_LUAN_ECG.tex (hai lan), hoac Tectonic.
\documentclass[12pt,a4paper,oneside]{report}
\usepackage{iftex}
\ifPDFTeX
  \usepackage[utf8]{inputenc}
  \usepackage[T5]{fontenc}
  \usepackage{lmodern}
\else
  \usepackage{fontspec}
  \IfFontExistsTF{Times New Roman}{\setmainfont{Times New Roman}}{\setmainfont{TeX Gyre Termes}}
  \IfFontExistsTF{Arial}{\setsansfont{Arial}}{\setsansfont{TeX Gyre Heros}}
  \IfFontExistsTF{Consolas}{\setmonofont[Scale=0.88]{Consolas}}{\setmonofont{Latin Modern Mono}}
\fi
\usepackage[vietnamese]{babel}
\usepackage[left=3cm,right=2cm,top=2.5cm,bottom=2.5cm]{geometry}
\usepackage{setspace}
\usepackage{amsmath,amssymb}
\usepackage{booktabs,longtable,tabularx,array}
\usepackage{graphicx,float}
\usepackage{xcolor}
\usepackage{enumitem}
\usepackage{tikz}
\usetikzlibrary{arrows.meta,positioning}
\usepackage{url}
\usepackage[unicode,hypertexnames=false,colorlinks=true,linkcolor=black,citecolor=blue!50!black,urlcolor=blue!50!black]{hyperref}
\hypersetup{pdftitle={Báo cáo khóa luận: Hệ thống ECG nhúng và đánh giá AF end-to-end},pdfauthor={Sinh viên thực hiện đề tài}}
\setstretch{1.15}
\setlength{\parindent}{1cm}
\setlength{\parskip}{0.25em}
\setlength{\tabcolsep}{4pt}
\renewcommand{\arraystretch}{1.15}
\setlist[itemize]{leftmargin=1.2cm,itemsep=0.15em,topsep=0.3em}
\setlist[enumerate]{leftmargin=1.2cm,itemsep=0.15em,topsep=0.3em}
\setcounter{secnumdepth}{2}
\setcounter{tocdepth}{1}
\emergencystretch=2em
\Urlmuskip=0mu plus 2mu
\newcolumntype{L}[1]{>{\raggedright\arraybackslash}p{#1}}
\newcommand{\artifact}[1]{\par\noindent\textit{Nguồn thực nghiệm:}\ \path{#1}\par}
\newcommand{\BA}{\operatorname{BA}}
\newcommand{\AF}{\operatorname{AF}}
\newcommand{\RR}{\operatorname{RR}}

\begin{document}
\begin{titlepage}
\centering
\vspace*{1cm}
{\Large BÁO CÁO KHÓA LUẬN\par}
\vspace{1.4cm}
{\LARGE\bfseries KHẢO SÁT VÀ XÂY DỰNG\par HỆ THỐNG PHÂN TÍCH ECG NHÚNG\par}
\vspace{0.5cm}
{\Large Phát hiện QRS, đánh giá chất lượng tín hiệu\par và sàng lọc rung nhĩ từ khoảng RR\par}
\vspace{1.3cm}
\begin{minipage}{0.86\textwidth}
\setstretch{1.2}
\textbf{Nội dung:} tổng hợp cơ sở lý thuyết, các phương pháp đã khảo sát, phần mềm và firmware đã xây dựng, kết quả kiểm chứng, những kết quả chưa đạt và hướng phát triển của đề tài.

\medskip
\textbf{Phạm vi cập nhật:} bằng chứng mã nguồn và artifact hiện có trong workspace đến ngày 31/10/2026. Đây là báo cáo thực hiện đề tài; những khối mới ở mức đề xuất được trình bày tách khỏi các kết quả đã đo.
\end{minipage}
\vfill
{\large Sinh viên thực hiện đề tài\par}
\vspace{0.4cm}
{\large Tháng 10 năm 2026\par}
\end{titlepage}

\pagenumbering{roman}
\chapter*{Tóm tắt}
\addcontentsline{toc}{chapter}{Tóm tắt}
Trong khóa luận này, em khảo sát và xây dựng một chuỗi xử lý điện tâm đồ (ECG) hướng tới thiết bị nhúng, từ thu nhận tín hiệu đến phát hiện đỉnh R, tính khoảng RR, ước lượng nhịp tim và sàng lọc rung nhĩ (AF). Trọng tâm thực nghiệm là đánh giá đồng thời độ chính xác, độ bao phủ, mức suy giảm khi có nhiễu và chi phí triển khai. Bên cạnh đó, em xây dựng cách đánh giá tổn hại do bộ khử nhiễu gây ra trên tín hiệu sạch và khảo sát khả năng giải thích mô hình bằng TreeSHAP.

Phần QRS đã có phép so sánh đầu vào thống nhất cho 26 phương pháp và số đo của 15 hiện thực C trên ESP32. Phần đánh giá chất lượng tín hiệu (SQI) bao gồm các quy tắc, đặc trưng hình thái, so khớp detector và mô hình học máy. Phần AF gồm các mô hình RR truyền thống, lượng tử hóa INT8, thử nghiệm ECGFounder và ConvNeXt một chuyển đạo. Các kết quả được giữ riêng theo đơn vị đánh giá: nhịp, cửa sổ hoặc giây trên toàn dòng thời gian.

Đợt lựa chọn detector cuối cùng gồm 296 job: 6 waveform AFDB với 19 điều kiện, 84 record LTAFDB để sàng lọc và 98 record SHDB-AF thuộc 93 bệnh nhân để xác nhận ứng viên đã khóa. Với Extra Trees INT8 và đầu ra xác nhận 3/4, detector firmware liên tục đạt BA lần lượt 0,8785; 0,9233; 0,8898. Không có ứng viên trong shortlist đáp ứng đầy đủ điều kiện thay thế đã khai báo. Vì vậy, cấu hình được giữ là detector firmware hiện tại kết hợp Extra Trees INT8; SQI được dùng như thông tin chất lượng bổ sung trong firmware, chưa được xem là gate AF đã được xác nhận.

Các kết quả trên là bằng chứng kỹ thuật từ replay và các phép đo kernel trên thiết bị. Đề tài chưa xác nhận hiệu quả lâm sàng, chưa chứng minh mọi khối server đã tích hợp, và chưa có phép đo mới của toàn bộ pipeline trên board tương ứng với đợt lựa chọn cuối cùng.

\tableofcontents
\listoftables
\listoffigures
\clearpage
\chapter*{Danh mục từ viết tắt}
\addcontentsline{toc}{chapter}{Danh mục từ viết tắt}
\begin{longtable}{L{2.4cm}L{12.4cm}}
\toprule
Từ viết tắt & Ý nghĩa sử dụng trong báo cáo\\\midrule
ECG & Điện tâm đồ (electrocardiogram).\\
AF & Rung nhĩ (atrial fibrillation); trong nhãn thử nghiệm, AF/AFIB là lớp dương.\\
AFL & Cuồng nhĩ (atrial flutter); được coi là lớp không AF trong các thử nghiệm AF-only.\\
QRS & Phức bộ biểu diễn hoạt động khử cực thất; đỉnh R là mốc thời gian dùng để tính RR.\\
RR, BPM & Khoảng giữa hai đỉnh R liên tiếp; số nhịp mỗi phút.\\
SQI & Chỉ số hoặc thuật toán đánh giá chất lượng tín hiệu.\\
BA, AUC & Độ chính xác cân bằng; diện tích dưới đường cong ROC.\\
PPV, F1 & Giá trị dự đoán dương; trung bình điều hòa giữa precision và recall.\\
MAE, RMSE & Sai số tuyệt đối trung bình; căn sai số bình phương trung bình.\\
PRD & Sai khác căn bình phương tương đối, biểu diễn theo phần trăm.\\
BW, MA, EM & Trôi đường nền, nhiễu cơ, nhiễu do chuyển động điện cực trong NSTDB.\\
AFE & Khối đầu vào tương tự; trong benchmark có mô phỏng một phần đặc tính đầu vào.\\
INT8, FP32 & Biểu diễn số nguyên 8 bit; số thực dấu phẩy động 32 bit.\\
XAI, SHAP & Giải thích mô hình; phương pháp quy đóng góp đặc trưng dựa trên giá trị Shapley.\\
Abstain & Trạng thái chưa đưa ra kết luận phân loại cho khoảng thời gian đang xét.\\
Host replay & Phát lại dữ liệu và chạy mã trên máy tính; khác với đo trực tiếp trên ESP32.\\
\bottomrule
\end{longtable}
\clearpage
\pagenumbering{arabic}

\chapter{Giới thiệu đề tài}
\section{Bài toán và động cơ thực hiện}
Thiết bị ECG một chuyển đạo có ưu điểm về kích thước và khả năng theo dõi liên tục, nhưng tín hiệu dễ thay đổi theo tiếp xúc điện cực, tư thế và chuyển động. Một đoạn tín hiệu có hình dạng tương đối dễ quan sát vẫn có thể chứa đỉnh R bị đặt sai vị trí; ngược lại, tín hiệu có biến thiên sinh lý mạnh không nhất thiết là nhiễu. Vì vậy, chỉ hiển thị BPM hoặc chỉ báo cáo một chỉ số phân loại AF chưa đủ để đánh giá một hệ thống hoàn chỉnh.

Mục tiêu của em là xây dựng một đường xử lý có thể kiểm chứng: mỗi quyết định AF phải truy ngược được tới ECG, các đỉnh R, cửa sổ RR, đặc trưng và phiên bản mô hình. Việc chọn mô hình cần xét cả trường hợp hệ thống chưa đủ dữ liệu để kết luận, thay vì loại các khoảng thời gian khó ra khỏi báo cáo. Hướng này phù hợp với yêu cầu của theo dõi ECG đeo người, nơi độ tin cậy phụ thuộc đồng thời vào chất lượng dữ liệu và thuật toán \cite{smital2020,daluwatte2016}.

\section{Mục tiêu và câu hỏi nghiên cứu}
Đề tài được tổ chức theo bốn câu hỏi:
\begin{enumerate}
\item Khi các detector nhận cùng tín hiệu và được chấm cùng quy tắc, phương pháp nào phù hợp với bài toán QRS/RR trên thiết bị?
\item Sai số của detector ảnh hưởng thế nào đến AF khi giữ nguyên mô hình và ngưỡng phân loại?
\item SQI giúp tăng độ tin cậy tới mức nào, và việc từ chối kết luận có làm bỏ sót thời gian AF hay không?
\item Chi phí Flash, RAM, thời gian tính toán và phạm vi hiện thực C có phù hợp với ESP32 đang sử dụng không?
\end{enumerate}
Đối với khử nhiễu, em bổ sung câu hỏi phụ: cải thiện SNR có đồng thời bảo toàn hình thái và khoảng RR hay không? Đối với XAI, mục tiêu là giải thích đúng quyết định của mô hình đã dùng, bao gồm cả các trường hợp dự đoán sai.

\section{Phạm vi và mức hoàn thành}
\begin{table}[H]
\centering\small
\caption{Phạm vi thực hiện và mức bằng chứng hiện có}
\begin{tabularx}{\textwidth}{L{3.1cm}X}
\toprule Hạng mục & Kết quả và phạm vi\\\midrule
Thu nhận và firmware & Có mã ADC, xử lý liên tục, AF INT8, SD/BLE và máy trạng thái; có smoke test và benchmark kernel trên ESP32.\\
Khảo sát QRS & Có 26 phương pháp ở phép so sánh common-input; 15 hiện thực C có số đo toàn cohort crop trên board.\\
AF downstream & Có so sánh cùng mô hình trên RR annotation/detector, ma trận nhiễu và phép xác nhận ứng viên theo bệnh nhân.\\
SQI & Có khảo sát, benchmark BUTQDB, kiểm tra port C và số đo tài nguyên; Ho đầy đủ chưa được port.\\
Khử nhiễu & Có giao thức CleanGuard/CleanGuard+, code baseline và smoke nhỏ; MECG-E mới là phương án được khảo sát.\\
Giải thích và server & TreeSHAP ngoại tuyến đã chạy; có prototype gateway BLE--SQLite. Chưa có bằng chứng dịch vụ server/XAI/fusion hai chiều hoàn chỉnh.\\
\bottomrule
\end{tabularx}
\end{table}

\section{Cách sử dụng bằng chứng trong báo cáo}
Em ưu tiên mã nguồn và các artifact đã lưu: protocol, hash, dự đoán, metric và audit. Các báo cáo cũ được dùng để mô tả quá trình phát triển; kết luận hiện tại dựa trên các lượt chạy đã sửa giao thức và kiểm tra lại. Cụ thể, bảng QRS có chuẩn hóa min--max khác miền đầu vào đã được thay bằng bảng common-input; kích thước image Extra Trees cũ bị gán nhầm đã được thay bằng số đo mới. Bản tổng hợp LaTeX ngày 10/10/2026 được xem là tài liệu lịch sử, không phải trạng thái triển khai hiện tại.

\chapter{Cơ sở lý thuyết và chỉ tiêu đánh giá}
\section{Từ ECG đến đỉnh R, RR và BPM}
ECG gồm các thành phần P, QRS và T. Trong phạm vi đề tài, em tập trung vào phức bộ QRS và timestamp đại diện của nhịp, chưa triển khai bộ phân định đầy đủ khởi đầu/kết thúc P--QRS--T. Detector QRS trả lời có một nhịp tại vùng tín hiệu nào; định vị R yêu cầu xác định chính xác mốc thời gian trong vùng đó. Hai yêu cầu này cần được đánh giá riêng.

Với các timestamp $t_i$ tính bằng ms:
\begin{equation}
RR_i=t_i-t_{i-1},\qquad BPM_i=\frac{60000}{RR_i}.
\end{equation}
Trong benchmark BPM hiện hành, một ước lượng được tính từ trung bình RR trong cửa sổ trượt 5 giây. Nếu mọi đỉnh R cùng lệch một lượng gần như nhau, RR và BPM vẫn có thể gần đúng. Vì vậy, BPM MAE thấp không chứng minh timestamp R chính xác tại dung sai $\pm50$ ms.

\section{Đánh giá QRS và sai số RR}
Các nhịp được ghép một-một với annotation trong dung sai thời gian đã khai báo. Mỗi nhịp chỉ được ghép tối đa một lần. Em sử dụng:
\begin{equation}
Se=\frac{TP}{TP+FN},\quad PPV=\frac{TP}{TP+FP},\quad
F_1=\frac{2TP}{2TP+FP+FN}.
\end{equation}
Dung sai $\pm50$ ms là tiêu chí chính về định vị; $\pm150$ ms là tiêu chí bổ sung. Không dùng annotation của tập kiểm tra để bù trễ riêng cho từng detector. RR MAE chỉ được tính trên các cặp nhịp liên tiếp khớp ở cả tham chiếu và dự đoán; vì đây là sai số có điều kiện, báo cáo luôn kèm tỷ lệ RR tham chiếu được bao phủ.

Với sai số $e_i$, các chỉ số MAE và RMSE là:
\begin{equation}
MAE=\frac{1}{N}\sum_{i=1}^{N}|e_i|,\qquad
RMSE=\sqrt{\frac{1}{N}\sum_{i=1}^{N}e_i^2}.
\end{equation}
Nếu một cửa sổ không có ước lượng BPM hợp lệ, cửa sổ đó vẫn nằm trong mẫu số của độ bao phủ; em không chuyển nó thành một dự đoán BPM bằng không để tính lỗi thông thường.

\section{Đặc trưng RR phục vụ phát hiện AF}
Mô hình nhúng sử dụng 30 RR và năm đặc trưng. Với $d_i=RR_{i+1}-RR_i$, pNN50 được biểu diễn dưới dạng tỷ lệ:
\begin{equation}
pNN50=\frac{1}{29}\sum_{i=1}^{29}\mathbf{1}(|d_i|>50\,\mathrm{ms}).
\end{equation}
Hai đại lượng Poincaré được tính từ phương sai mẫu:
\begin{equation}
SD_1=\sqrt{\tfrac12\operatorname{Var}(d)},\qquad
SD_2=\sqrt{\max(0,2\operatorname{Var}(RR)-SD_1^2)}.
\end{equation}
Các đặc trưng còn lại là $SD_1/SD_2$, median RR và Shannon entropy. Khi $SD_2=0$, tỷ lệ được quy ước bằng không. Runtime dùng năm bin trải từ RR nhỏ nhất đến lớn nhất để tính:
\begin{equation}
H=-\sum_{b:p_b>0}p_b\log_2p_b.
\end{equation}
Những đặc trưng này mô tả mức bất quy tắc của nhịp. Chúng không quan sát trực tiếp sóng P, do đó biến thiên RR của các nhịp không AF cũng có thể gây dương tính giả. Công thức dùng float32 trong firmware; kiểm tra số học cần giữ đúng cách tích lũy, đặc biệt khi $SD_2$ là hiệu của các đại lượng gần nhau.

\section{Abstain, độ bao phủ và recall AF hiệu dụng}
Em dùng từ \emph{abstain} cho trạng thái hệ thống chưa kết luận AF hay không AF. Các nguyên nhân có thể là chưa đủ RR, RR không hợp lệ làm reset hoặc gate chất lượng từ chối đầu ra. Trạng thái này khác một dự đoán âm tính.

\begin{samepage}
Gọi $\mathcal{T}$ là tập giây có nhãn, $\mathcal{A}$ là phần có dự đoán được chấp nhận và $y_t$ là nhãn:
\begin{align}
Coverage&=\frac{|\mathcal{A}|}{|\mathcal{T}|},\\
AF\ retention&=\frac{|\{t\in\mathcal{A}:y_t=1\}|}{|\{t\in\mathcal{T}:y_t=1\}|},\\
Recall_{\mathrm{AF,eff}}&=\frac{TP}{|\{t\in\mathcal{T}:y_t=1\}|}.
\end{align}
\end{samepage}
BA được tính trên phần có kết luận:
\begin{equation}
BA=\frac{1}{2}\left(\frac{TP}{TP+FN}+\frac{TN}{TN+FP}\right).
\end{equation}
Ví dụ, trong 100 giây AF, hệ thống bắt đúng 80 giây, dự đoán âm tính 10 giây và abstain 10 giây. Recall có điều kiện bằng $80/90$, trong khi recall AF hiệu dụng bằng $80/100$. Báo cáo chỉ có recall có điều kiện sẽ che mất một phần thời gian không được phục vụ.

\section{Xác nhận theo thời gian và metric episode}
Mỗi lần có thêm 15 RR hợp lệ, mô hình được cập nhật trên 30 RR gần nhất. Đầu ra xác nhận chỉ bật khi đã có đủ bốn dự đoán và ít nhất ba dự đoán là AF. Ví dụ, chuỗi AF--không AF--AF--AF đủ ba phiếu; chuỗi AF--không AF--AF--không AF chưa đủ. Con số 3/4 là ngưỡng số phiếu để bật cảnh báo; điểm dự đoán của từng cửa sổ được lưu riêng.

Bốn lần suy luận đầu tiên xuất hiện sau 30, 45, 60 và 75 RR hợp lệ. Các cửa sổ này chồng lấn, vì vậy bốn dự đoán không phải bốn quan sát độc lập. Sau startup, lần xác nhận sớm nhất cần 75 RR hợp lệ; khi RR không hợp lệ, lịch sử xác nhận bị xóa. Quy tắc làm ổn định đầu ra nhưng có thể làm tăng độ trễ hoặc bỏ qua episode ngắn, nên cần được đánh giá cùng metric episode.

Đánh giá theo episode bổ sung tỷ lệ episode AF bắt được, độ trễ bắt đầu và số lần bắt đầu cảnh báo sai trên một giờ không AF. Độ trễ chỉ tính trên các episode thực sự được bắt; episode bỏ sót vẫn ở mẫu số của recall episode. Episode đã tồn tại tại đầu record được giữ và ghi nhận giới hạn quan sát đầu vào. Lưới 1 giây giới hạn độ phân giải của các số đo này.

\chapter{Khảo sát các hướng tiếp cận}
\section{Detector QRS cổ điển và hiện thực nhúng}
Pan--Tompkins sử dụng lọc, đạo hàm, năng lượng và ngưỡng thích nghi để phát hiện QRS \cite{pan1985}. Các phương pháp theo Elgendi sử dụng các cửa sổ trung bình trượt và quy tắc dựa trên kiến thức về phức bộ QRS \cite{elgendi2013}. Đề tài không coi mọi hiện thực cùng tên là tương đương: bản Pan/Elgendi của thư viện \texttt{py-ecg-detectors}, bản trong WIM và detector firmware là các nhánh khác nhau, có thể khác cả bộ lọc lẫn quy tắc đặt timestamp.

Em khảo sát thêm WFDB GQRS/XQRS, FastQRS 0.3.0 và Pan--Tompkins++. Các detector batch thuận tiện cho so sánh nhưng có thể dùng lọc hai chiều; kết quả của một khối chỉ sẵn có sau khi khối được thu xong. Khi đánh giá AF downstream, em ghi nhận độ trễ sẵn có này thay vì gán kết quả về đầu khối.

\begin{longtable}{L{3.1cm}L{6.7cm}L{4.9cm}}
\caption{Các nhóm detector đã khảo sát và đánh giá}\label{tab:surveyqrs}\\
\toprule Nhóm & Phương pháp trong common-input & Phạm vi kết luận\\\midrule\endfirsthead
\toprule Nhóm & Phương pháp trong common-input & Phạm vi kết luận\\\midrule\endhead
Firmware và thư viện & Production C; Pan--Tompkins/Elgendi từ py-ecg-detectors; WFDB XQRS/GQRS & Baseline thực tế và tham chiếu phần mềm; cần giữ đúng lead, fs và miền đơn vị.\\
Detector gọn & FastQRS 0.3.0; Pan--Tompkins++; WIM Nguyen, Elgendi, Hamilton, Arteaga--Falconi, Pan--Tompkins & Có C port cho nhiều phương pháp; chất lượng timestamp và chi phí bộ nhớ không giống nhau.\\
Biến đổi và quy tắc & WIM Park, Kumari, Shaik, Xia, Xu, Zhai & Bao gồm các phép biến đổi, lọc và heuristic khác nhau; điểm ở dung sai rộng có thể che lỗi định vị.\\
Học sâu & ECG2RR LSTM; SWT--SepConv; WIM Laitala, Xiang, Zahid, Celik, Han CNN, Han RNN & Có kết quả host; khả năng chạy MCU phải kiểm tra riêng tensor, operator và bộ nhớ.\\
\bottomrule
\end{longtable}

\section{Mô hình QRS học sâu và fusion}
ECG2RR/Laitala sử dụng LSTM cho phát hiện R \cite{laitala2020}. SWT--SepConv kết hợp stationary wavelet transform với mạng convolution tách được \cite{yun2022}. Em dùng các mô hình này để khảo sát mức chất lượng có thể đạt trên dữ liệu và để đối chiếu sai số của phương pháp nhẹ. Checkpoint SWT đã được huấn luyện với MITDB, INCART và QTDB; điểm cao trên các bộ dữ liệu này không được gọi là tổng quát hóa trên dữ liệu chưa thấy.

Đề tài cũng thử kết hợp detector bằng voting, cặp detector và cross-attention. Ý tưởng là tận dụng sự bổ sung về vị trí peak, nhưng nhiều detector có thể cùng mắc một loại lỗi; lấy hợp tất cả peak sẽ làm tăng dương tính giả. Do đó, fusion phải được so sánh với detector đơn mạnh nhất trên cùng dữ liệu và split, thay vì chỉ so với một baseline yếu.

\section{SQI theo quy tắc, hình thái và độ đồng thuận}
Em khảo sát ba nhóm chính. Nhóm quy tắc sử dụng giới hạn nhịp/RR và tính hợp lý của tín hiệu. Nhóm hình thái dùng tương quan template, kurtosis, skewness hoặc phân bố năng lượng phổ. Nhóm đồng thuận so sánh kết quả của các detector độc lập.

Orphanidou kết hợp các quy tắc và mức tương đồng hình thái \cite{orphanidou2015}. Clifford khảo sát các SQI và data fusion để đánh giá mức chấp nhận ECG \cite{clifford2012}. Zhao sử dụng heuristic fusion và đánh giá mờ cho ECG một chuyển đạo \cite{zhao2018}. Ho đề xuất khai thác mức nhất quán của RR trong ECG telehealth \cite{ho2025}. Trong workspace, phép chạy Ho tham chiếu sử dụng NeuroKit 0.2.13 với cặp detector UNSW và NeuroKit; ngưỡng ít nhất 80\% sample tốt là policy của dự án.

Điểm cần thận trọng là biến thiên do loạn nhịp có thể bị nhầm với chất lượng kém \cite{daluwatte2016}. Vì vậy, một SQI có BA tốt trên nhãn chất lượng chưa tự chứng minh rằng gate của nó sẽ giữ được thời gian AF. Tác động downstream được đánh giá riêng ở Chương~\ref{chap:af}.

\section{Mô hình AF và khả năng giải thích}
Đối với đầu vào RR, em so sánh Logistic Regression, Random Forest, Extra Trees và XGBoost. Extra Trees tạo ensemble cây ngẫu nhiên hóa mạnh \cite{geurts2006}; mô hình có thể chuyển thành các bảng hằng số và phép so sánh phù hợp cho C. Logistic Regression có ưu thế kích thước và diễn giải tuyến tính; các mô hình cây cho phép khảo sát tương tác phi tuyến.

Đối với waveform một chuyển đạo, em khảo sát ECGFounder và kiến trúc ConvNeXt lấy cảm hứng từ Kraft và Rumm \cite{ecgfounder2025,kraft2025}. Thử nghiệm ECGFounder đóng băng encoder và huấn luyện đầu logistic; không phải fine-tune toàn mạng. ConvNeXt trong dự án là hiện thực lấy cảm hứng từ bài báo; không được trình bày như tái hiện nguyên mã hoặc checkpoint gốc. Bài Kraft và Rumm gộp AF/AFL, trong khi nhiệm vụ của dự án chỉ dùng AF là dương tính; khác biệt nhãn này ảnh hưởng khả năng so sánh.

Đối với giải thích, SHAP phân bổ đóng góp của các đặc trưng cho đầu ra mô hình \cite{shap2017}; TreeSHAP khai thác cấu trúc cây để tính hiệu quả \cite{treeshap2020}. Giải thích mô hình cho biết yếu tố nào làm điểm dự đoán tăng hoặc giảm; nó không xác nhận cơ chế sinh lý gây ra bệnh.

\section{Khử nhiễu và định hướng phân định waveform}
Các baseline FIR và wavelet được khảo sát nhằm cân bằng loại nhiễu và bảo toàn tín hiệu. MECG-E là mô hình Mamba dành cho loại trôi đường nền \cite{mecge2024}; đề tài mới khảo sát khả năng đưa mô hình này vào nhánh server, chưa có checkpoint hay metric cục bộ để kết luận hiệu quả trên dữ liệu đang dùng.

Hướng SQI gắn với phân định P--QRS--T được khảo sát qua tài liệu của Xie và các bộ delineation \cite{xie2021,martinez2004}. Trong trạng thái hiện tại, em chưa triển khai đầy đủ delineation hay chứng minh các mức chất lượng Q2/Q1/Q0 dẫn tới metric fiducial tương ứng. Các nội dung này thuộc định hướng mở rộng, không được gộp vào kết quả QRS đã chạy.

Em cũng khảo sát ECGdeli, một toolbox MATLAB phân định các thành phần ECG, và nghiên cứu so sánh delineation trên dữ liệu smartwatch \cite{ecgdeli2021,jaeger2024}. Các tài liệu này gợi ý cần chấm riêng sai số mốc bắt đầu, đỉnh và kết thúc của từng sóng, đồng thời đánh giá ảnh hưởng của chuyển đạo và chất lượng tín hiệu. Đề tài chưa chạy hai nhánh đó trên dữ liệu thiết bị; chúng được dùng để xây dựng kế hoạch đánh giá sau này.
\artifact{docs/SQA_BASED_ECG_DELINEATION_PIPELINE.md}

\chapter{Dữ liệu và phương pháp thực nghiệm}
\section{Các bộ dữ liệu đã sử dụng}
\begin{longtable}{L{2.1cm}L{3.8cm}L{8.6cm}}
\caption{Dữ liệu và vai trò trong các thí nghiệm chính}\label{tab:datasets}\\
\toprule Dataset & Quy mô dùng trong báo cáo & Vai trò và giới hạn\\\midrule\endfirsthead
\toprule Dataset & Quy mô dùng trong báo cáo & Vai trò và giới hạn\\\midrule\endhead
MITDB & 46 record MLII, 45 nhóm bệnh nhân & QRS/RR/BPM. Record 201/202 cùng bệnh nhân; chọn MLII theo header, không giả định channel 0 \cite{mitdb}.\\
QTDB & 13 đoạn MLII ở phép so sánh hiện hành; các nhánh lịch sử có cohort khác & Cohort MLII này có nguồn MITDB, dùng kiểm tra hồi quy; không cộng như external độc lập \cite{qtdb}.\\
NSTDB & BW, MA, EM; sáu mức SNR & Nhiễu cục bộ có kiểm soát; không thay thế nhiễu vận động thực tế của thiết bị \cite{nstdb}.\\
AFDB & 25 record annotation cho tập feature; 6 waveform test hiện có, 60,393 giờ & Split mô hình 18 train/7 test; record 00735 thiếu waveform. Nhãn rhythm dùng .atr; .qrs là mốc nhịp tự động chưa hiệu đính toàn bộ \cite{afdb}.\\
LTAFDB & 84 record, 1.960,6 giờ, 128 Hz & Đánh giá AF ngoài AFDB và sàng lọc detector. Chưa có ánh xạ bệnh nhân cục bộ, nên bootstrap theo record \cite{ltafdb}.\\
SHDB-AF & 98 record có annotation, 93 bệnh nhân, 2.330,7 giờ, 200 Hz & Đánh giá ngoài và xác nhận một ứng viên đã khóa; dùng mapping từ AdditionalData.csv \cite{shdbaf}.\\
BUTQDB & 18 record/15 người, 478,73 giờ ECG; 35.773 cửa sổ có consensus (99,37 giờ) & SQI theo nhãn chất lượng, split theo người. Trong lượt này không có nhãn AF/R-peak độc lập để chấm AF/BPM \cite{butqdb}.\\
CinC2011 & 998 record Set A, 10 giây, 500 Hz & Nhãn toàn bản ghi 12 lead. Dùng Lead II để kiểm tra port không tạo ra nhãn SQI riêng cho Lead II \cite{cinc2011}.\\
ECG tự thu & Các CSV của thiết bị/AD8232, nhiều cohort xử lý & Dùng kiểm tra luồng dữ liệu và đầu ra SQI. Chưa có nhãn chất lượng, R-peak và AF độc lập để ước lượng accuracy.\\
\bottomrule
\end{longtable}

Các cohort BUT PPG và VitalDB chưa có đủ waveform/annotation cần thiết trong inventory đã kiểm tra; em không biến các phần tải chưa đủ thành một external benchmark. Số file có sẵn và số record đã được đánh giá là hai trạng thái riêng.

\section{Đầu vào chung và ma trận nhiễu}
Phép QRS common-input đưa tín hiệu về 360 Hz, lọc high-pass nhân quả 0,5 Hz, mô phỏng 200 counts/mV với baseline 2048 và clip 12 bit. Đó là mô phỏng một phần AFE, không phải mô hình đã hiệu chuẩn của toàn mạch vật lý. Các phương pháp nhận cùng mV quy đổi từ ADC; các biến đổi nội tại của từng thuật toán vẫn được giữ.

Ở đợt lựa chọn AF downstream cuối cùng, em đưa cùng đầu ra notch firmware cho các ứng viên batch. Detector production giữ trạng thái liên tục và dùng chính notch nội tại đó. Không thêm một tầng khử nhiễu hoặc SQI gate vào phép so sánh này.

Nhiễu được trộn thành burst 10 giây, lặp mỗi 30 giây và bắt đầu từ giây thứ 10. Mức SNR là 24, 18, 12, 6, 0 và $-6$ dB cho từng nguồn BW/MA/EM; biên độ dựa trên RMS cục bộ sau khi trừ trung bình. Vị trí nhiễu được xác định bằng hash record để các phương pháp nhận cùng hiện thực nhiễu. Em báo riêng toàn record, phần burst và phần không thêm nhiễu; không dùng BA tăng khi có nhiễu như bằng chứng mô hình tốt lên.

\section{Tách các giao thức đánh giá}
\begin{table}[H]
\centering\small
\caption{Các giao thức được giữ riêng khi diễn giải kết quả}
\begin{tabularx}{\textwidth}{L{3.2cm}X}
\toprule Giao thức & Đơn vị đánh giá\\\midrule
QRS crop common-input & 59 đoạn 10 giây tại giây 30--40; ghép nhịp tại dung sai đã khóa.\\
QRS toàn record & Dòng tín hiệu đầy đủ, reset một lần mỗi record, giữ startup.\\
AF feature benchmark & Cửa sổ 30 RR được suy từ annotation, nhãn thuần theo cửa sổ; 20.194 cửa sổ test.\\
AF timeline & Lưới 1 giây, lấy kết quả mới nhất đã sẵn có theo processing time; giữ giây chuyển nhịp.\\
AF waveform học sâu & Cửa sổ ECG thuần nhãn dài 10 giây; loại unknown và cửa sổ cắt qua chuyển nhịp.\\
SQI BUTQDB & Cửa sổ 10 giây có consensus; chia nhóm theo bệnh nhân, target chất lượng phục vụ QRS.\\
ESP32 kernel & Thời gian thực thi detector/scorer/classifier trên image benchmark và input đã quy định.\\
\bottomrule
\end{tabularx}
\end{table}

\section{Split, khóa tham số và kiểm chứng}
Mô hình AF được huấn luyện trên 18 record và kiểm tra trên 7 record của AFDB; chỉ 6 record test có waveform cục bộ. Tham số mô hình, feature và ngưỡng được giữ cố định khi thay nguồn RR. Dữ liệu AFDB đã được xem nhiều lần trong phát triển, vì vậy em không gọi đây là kiểm định tiến cứu.

Đợt detector-selection sử dụng AFDB và LTAFDB ở bước sàng lọc. Sau đó, một ứng viên được khóa trước khi chạy bước SHDB-AF. Các dataset ngoài từng được dự án đánh giá với phương pháp khác; SHDB-AF được giữ ngoài \emph{bước chọn detector của lượt này}, không phải tập dữ liệu chưa từng được dự án tiếp cận. Bootstrap LTAFDB lấy mẫu theo record; SHDB-AF lấy mẫu ghép cặp theo 93 bệnh nhân để giữ các record cùng người đi cùng nhau.

Mỗi lượt chạy lưu protocol, hash source/input, đầu ra và audit. Kiểm chứng mã C gồm so với công thức độc lập, replay lại trạng thái, kiểm tra số liệu gộp từ từng record và xác nhận timestamp không sử dụng dự đoán tương lai. Các repository bên thứ ba được giữ riêng trong thư mục external; em sử dụng adapter thay vì sửa trực tiếp upstream.

\chapter{Xây dựng hệ thống và firmware}
\section{Phần cứng và luồng xử lý}
Hệ thống hiện có cấu hình đầu vào AD8232 và ESP32. Source ADC dùng ADC1/channel 0, độ phân giải 12 bit và cấu hình suy hao 12 dB; các comment về chân nối còn mơ hồ nên báo cáo không suy ra số GPIO chỉ từ một tên biến. Tần số xử lý mục tiêu là 360 Hz. Tín hiệu được trừ baseline 2048, qua notch 49--51 Hz và đưa vào nhánh QRS.

Detector production là một hiện thực nhân quả theo hướng Pan-style: lọc dải phục vụ QRS, năng lượng, các cửa sổ trung bình trượt 100/600 ms, ngưỡng thích nghi và căn R trên lịch sử tín hiệu. Cấu hình có refractory 300 ms, vùng tối thiểu 60 ms và lịch sử raw 128 mẫu. Đây là hiện thực dự án; không tự nhận là tái hiện nguyên vẹn thuật toán Pan--Tompkins gốc.

\begin{figure}[H]
\centering
\begin{tikzpicture}[box/.style={draw,rounded corners,align=center,text width=3.1cm,minimum height=1cm,font=\small},arr/.style={-{Latex[length=2.5mm]},thick}]
\node[box] (adc) at (0,0) {Thu ECG\\ADC 360 Hz};
\node[box] (pre) at (4.1,0) {Trừ baseline\\Notch 49--51 Hz};
\node[box] (qrs) at (8.2,0) {Detector QRS\\firmware liên tục};
\node[box] (rr) at (8.2,-1.9) {Đỉnh R\\RR / BPM};
\node[box] (feat) at (4.1,-1.9) {30 RR\\5 đặc trưng};
\node[box] (et) at (0,-1.9) {Extra Trees INT8\\ngưỡng đã khóa};
\node[box] (confirm) at (0,-4.3) {4 dự đoán gần nhất\\ít nhất 3 phiếu AF};
\node[box] (result) at (4.1,-4.3) {Kết quả / độ hợp lệ\\timestamp / event};
\node[box] (io) at (8.2,-4.3) {Ghi SD\\gửi BLE};
\draw[arr] (adc)--(pre);\draw[arr] (pre)--(qrs);\draw[arr] (qrs)--(rr);
\draw[arr] (rr)--(feat);\draw[arr] (feat)--(et);\draw[arr] (et)--(confirm);
\draw[arr] (confirm)--(result);\draw[arr] (result)--(io);
\end{tikzpicture}
\caption{Đường xử lý chính được giữ sau phép lựa chọn downstream AF. SQI chạy như nhánh telemetry độc lập và chưa là gate chặn AF của cấu hình đã đánh giá.}
\label{fig:pipeline}
\end{figure}

\section{Tích hợp AF và trạng thái hợp lệ}
Sau mỗi đỉnh R, runtime cập nhật RR. Khoảng ngoài 250--2500 ms làm reset hàng đợi và lịch sử xác nhận. Khi đủ 30 RR và đến stride 15, runtime tính năm feature, lượng tử hóa và chạy Extra Trees INT8. Mô hình gồm 80 cây, 17.712 node; score nguyên được so với ngưỡng 10.240. Chuẩn hóa score theo $80\times255$ chỉ giúp biểu diễn đầu ra; ngưỡng INT8 là cấu hình đã khóa, không phải ngưỡng float được chuyển đổi tùy ý sau kiểm tra.

Telemetry phân biệt dự đoán thô, trạng thái xác nhận và trạng thái đủ dữ liệu. Khi chưa đủ RR, hệ thống phải hiển thị chưa đủ dữ liệu thay vì âm tính. Source hiện chưa có timeout độc lập khi không có peak mới; vì vậy một kết quả valid có thể cũ. Dữ liệu lưu kèm timestamp và audit tuổi đầu ra là cần thiết. Mốc 30 giây trong phân tích độ nhạy ở chương sau chưa phải TTL đã triển khai.

\section{SQI shadow trong source hiện tại}
SQI firmware hiện là Orphanidou-style, nhận tín hiệu đã notch và peak firmware. Source hiện dùng cửa sổ 10 giây với hop 1 giây, trong khi một số README lịch sử mô tả cửa sổ không chồng lấn. Báo cáo lấy cấu hình trong source làm hiện trạng.

AF được cập nhật độc lập trước SQI. SQI cung cấp thông tin unusable/usable/insufficient và các đại lượng chất lượng; nó không reset hoặc chặn nhánh AF production. Điều này khác các thí nghiệm \emph{output gate} trên host, nơi em chủ động chuyển một số giây thành abstain để đo tác động.

\section{Phân chia tác vụ, SD và BLE}
Core 1 thực hiện thu nhận và xử lý ECG; Core 0 phụ trách điều khiển, ghi SD và BLE. Máy trạng thái có INIT/IDLE/RECORD; hệ thống vào RECORD khi có START bằng nút hoặc BLE, thay vì tự ghi ngay khi boot. Khi STOP, hàng đợi được drain trước khi unmount.

Nhánh SD lưu full-rate, gom 32 dòng, dùng tên file 8.3 và chia segment 60 giây. Có event journal và cơ chế giữ dữ liệu lịch sử để đưa dấu R về đúng timestamp. Nhánh BLE v1 dùng gói notification 20 byte: gói waveform chứa năm mẫu Int16 với sequence/timestamp, còn gói event mang BPM/RR/SQI/AF. BLE ưu tiên event và không được mô tả như kênh truyền full-rate 360 Hz trong mọi cấu hình.

Prototype gateway có khả năng giải mã gói, ghi SQLite và phát hiện sequence gap. Một client BLE legacy 16 byte vẫn còn trong workspace; giao thức này không được gộp với v1 20 byte thành cùng một hệ thống đã tích hợp. Các bài kiểm tra SD/BLE ngắn xác nhận đường dữ liệu nhưng chưa tương đương endurance test 24 giờ.

\section{Thiết kế server dự kiến}
Thiết kế mong muốn gồm nhận ECG gốc với timestamp, phân tích waveform bổ sung và trả kết quả có phiên bản mô hình. Nếu dùng ECGFounder, đầu vào phải là waveform; Extra Trees nhận đặc trưng RR riêng. TreeSHAP cho Extra Trees cần giữ cả đầu ra float và INT8 khi hai mô hình không đồng thuận.

Trong trạng thái hiện tại, em có mô hình/offline explanation và gateway prototype, nhưng chưa có bằng chứng dịch vụ nhận ECG, khử nhiễu MECG-E, suy luận hai mô hình, fusion và trả giải thích đã chạy end-to-end. Các khối fusion/server trên sơ đồ Draw.io vì vậy được hiểu là thiết kế mở rộng.

\chapter{Thực nghiệm khử nhiễu và bảo toàn tín hiệu}
\section{Giao thức CleanGuard và CleanGuard+}
Với tín hiệu tham chiếu $x$, nhiễu $n$ và bộ khử nhiễu $D$, em đánh giá cả $D(x+n)$ lẫn $D(x)$. Nhánh thứ nhất đo lợi ích khử nhiễu; nhánh thứ hai đo tổn hại trên tín hiệu sạch. Các đại lượng cơ bản là:
\begin{align}
SNR_{out}&=10\log_{10}\frac{\sum x_i^2}{\sum(x_i-D(x+n)_i)^2},\\
PRD&=100\sqrt{\frac{\sum(x_i-\widehat{x}_i)^2}{\sum x_i^2}}.
\end{align}
Code còn báo RMSE, sai số đạo hàm và cosine similarity. CleanGuard được giữ như một vector metric, chưa bị gộp thành một điểm duy nhất với trọng số chọn sau khi xem kết quả.

CleanGuard+ bổ sung RMSE xấu nhất trên cửa sổ cục bộ 1 giây với bước 0,5 giây, PRD cục bộ lớn nhất và RMSE ở giây đầu/cuối. Mục đích là phát hiện lỗi ngắn hoặc biên overlap-add có thể bị trung bình toàn cửa sổ che khuất. Chỉ số hiệu chỉnh dịch tối đa 20 ms là chẩn đoán phụ; metric chính giữ căn mẫu gốc, không dùng DTW.

Sau xử lý, cùng detector Pan--Tompkins được chạy trên $D(x+n)$ và $D(x)$ để đo QRS/RR. Cách này cho phép phân biệt cải thiện waveform với cải thiện thật ở tác vụ sau đó.

\section{Smoke test đã có}
Artifact hiện có gồm bốn cửa sổ 10 giây từ QTDB sel123 và sel233, hai cửa sổ mỗi record, thêm MA ở 6 dB và ghép nhịp tại $\pm150$ ms. Các cửa sổ thực tế dùng annotation .atr. Runner legacy ưu tiên .pu1, rồi .pu0, sau đó mới .atr; các annotation .pu1/.pu0 do ecgpuwave tạo tự động. Khi mở rộng, phải khóa nguồn tham chiếu rõ ràng và không tự coi các annotation tự động là nhãn chuyên gia. Các số trong bảng là trung bình theo bốn cửa sổ; RR MAE được tính ở nhánh tín hiệu có nhiễu sau xử lý.

\begin{table}[H]
\centering\small
\caption{Smoke khử nhiễu: bốn cửa sổ QTDB, MA 6 dB}
\begin{tabular}{lrrrr}
\toprule Phương pháp & $SNR_{out}$ (dB) & Gain (dB) & PRD sạch (\%) & RR MAE (ms)\\\midrule
Không khử nhiễu & 6,000 & 0,000 & 0,000 & 11,813\\
FIR HP 0,5 Hz & 6,216 & 0,216 & 43,082 & 11,813\\
FIR 0,5--40 Hz & 6,333 & 0,333 & 43,366 & 10,938\\
Wavelet db4 & 6,437 & 0,437 & 3,934 & 7,875\\
Wavelet sym4 & 6,439 & 0,439 & 3,894 & 6,500\\
sym4 BW-tuned & 7,542 & 1,542 & 8,718 & 11,813\\
sym9 cg15 & 8,407 & 2,407 & 13,917 & 11,938\\
sym5 cg20 & 8,279 & 2,279 & 19,783 & 11,938\\
db4 SNR & 8,107 & 2,107 & 25,604 & 11,938\\
\bottomrule
\end{tabular}
\end{table}

Gain SNR tăng không luôn đi cùng RR MAE giảm. Chẳng hạn, sym9 cg15 đạt gain 2,407 dB nhưng RR MAE lớn hơn sym4. PRD cục bộ lớn nhất của FIR HP khoảng 76,568\%, trong khi db4 khoảng 10,610\%. Những số này minh họa lý do cần đo tổn hại sạch; chúng chưa đủ để khẳng định phương pháp tốt nhất trên QTDB. Mọi phương pháp đều đạt F1 bằng 1 ở dung sai rộng trong smoke nhỏ này.
\artifact{experiment/report/outputs/rpeak_baseline_smoke/summary.csv}

\section{Những phần chưa hoàn tất}
Em đã có runner và baseline nhưng chưa tìm thấy output đầy đủ của benchmark CleanGuard diện rộng tương ứng trong workspace. Chưa có bằng chứng MECG-E được chạy trên waveform cục bộ và nối vào AF. Đóng góp hiện tại của phần này là giao thức, metric và minh họa có kiểm soát; chưa phải một phương pháp khử nhiễu mới đã được chứng minh ưu thế.

\chapter{Kết quả phát hiện QRS, RR và BPM}
\section{So sánh common-input trên các đoạn 10 giây}
Phép chạy gồm 59 record-crop: 46 MITDB MLII và 13 QTDB MLII, mỗi đoạn từ giây 30 đến 40. Tổng số hàng đã kiểm tra là 1.534 cho 26 phương pháp. Bảng~\ref{tab:qrscommon} trình bày một số phương pháp đại diện; danh sách khảo sát đầy đủ ở Bảng~\ref{tab:surveyqrs}.

\begin{table}[H]
\centering\small
\caption{MITDB crop common-input: không đồng nhất với kết quả toàn record}\label{tab:qrscommon}
\begin{tabular}{lrrrr}
\toprule Phương pháp & F1 50 ms & F1 150 ms & BPM MAE & Host p95 (ms)\\\midrule
Production C & 0,9371 & 0,9580 & 0,835 & 0,367\\
WFDB XQRS & 0,9810 & 0,9897 & 0,606 & 18,937\\
ECG2RR LSTM & 0,9760 & 0,9777 & 1,120 & 300,011\\
FastQRS & 0,9914 & 0,9932 & 0,927 & 2,855\\
Pan--Tompkins++ & 0,9780 & 0,9865 & 1,695 & 2,083\\
SWT--SepConv & 0,9932 & 0,9949 & 0,438 & 23,348\\
WIM Nguyen & 0,9846 & 0,9966 & 0,706 & 2,137\\
WIM Zhai & 0,9889 & 0,9974 & 0,539 & 2,805\\
WIM Laitala & 0,9762 & 0,9966 & 0,551 & 199,986\\
\bottomrule
\end{tabular}
\end{table}

Thời gian host là thời gian gọi detector trên máy tính cho một mảng 10 giây sau warm-up, không phải latency ESP32. Điểm F1 nghiêm ngặt cho thấy nhiều ứng viên định vị R tốt hơn production trên crop, nhưng đây chưa phải căn cứ chọn detector cho AF. QTDB ở phép chạy này là một cohort hồi quy có nguồn MITDB, nên điểm gần hoàn hảo ở QTDB không tạo ra một external test độc lập.

Một số hiện thực WIM như Shaik/Xu có chênh lệch lớn giữa dung sai 50 và 150 ms; hiện tượng này phản ánh lỗi đặt mốc, dù detector vẫn tìm vùng QRS. Các adapter Park/Zahid/Kumari đã được kiểm tra riêng cùng đầu vào; thay đổi tương thích API không được xem là huấn luyện lại mô hình.

\section{Kết quả toàn record}
\begin{table}[H]
\centering\small
\caption{MITDB toàn record: giữ startup và reset một lần mỗi record}
\begin{tabular}{lrrrrr}
\toprule Phương pháp & F1 50 ms & F1 150 ms & RR MAE & Bao phủ RR & BPM MAE\\\midrule
Production C & 0,96045 & 0,98775 & 2,915 & 93,46\% & 1,798\\
XQRS & 0,98088 & 0,99256 & 2,941 & 96,21\% & 0,761\\
GQRS & 0,97316 & 0,99523 & 2,941 & 95,66\% & 0,633\\
ECG2RR & 0,97294 & 0,97694 & 2,828 & 93,00\% & 2,311\\
FastQRS & 0,98942 & 0,99249 & 4,807 & 97,03\% & 1,073\\
Pan--Tompkins++ & 0,88376 & 0,89128 & 3,548 & 74,67\% & 17,680\\
SWT--SepConv & 0,99825 & 0,99849 & 1,902 & 99,88\% & 0,098\\
\bottomrule
\end{tabular}
\end{table}
RR MAE ở bảng tính bằng ms trên cặp liên tiếp khớp. Kết quả FastQRS cho thấy F1 cao không bảo đảm RR MAE nhỏ nhất. Sự thay đổi của Pan--Tompkins++ giữa crop và toàn record cũng cho thấy không thể chọn một số tốt nhất từ các giao thức khác nhau rồi xếp cùng bảng.

Các nhánh phát triển trước gồm Adaptive PT v1/v2 và bản portable int16/float32. Golden-vector parity giúp xác nhận chuyển mã trên bộ vector đã thử, nhưng bản portable chưa vượt đủ điều kiện chất lượng để thay production. Các kết quả lịch sử này được giữ để giải thích quá trình phát triển, không thay thế phép lựa chọn downstream cuối cùng.

\section{Phép thử fusion và kết quả âm}
Em đã thực hiện phép 5-fold theo 45 nhóm bệnh nhân trên 59 crop, dùng các dự đoán lưu sẵn của 23 detector. Cross-attention nhận vote và độ lệch thời gian của peak, sau đó chấp nhận/từ chối candidate; timestamp cuối là median candidate. Đây là fusion peak, không phải kết hợp trực tiếp các giá trị BPM.

\begin{table}[H]
\centering\small
\caption{Fusion exploratory OOF trên cùng 59 crop}
\begin{tabular}{lrr}
\toprule Phương pháp & F1 $\pm50$ ms & BPM MAE\\\midrule
SWT--SepConv & 0,9934 & 0,407\\
FastQRS + SWT median & 0,9926 & 0,687\\
WIM Zhai & 0,9913 & 0,437\\
FastQRS & 0,9906 & 0,783\\
Cross-attention 23 detector & 0,9657 & 4,508\\
Vote ít nhất 5 & 0,9005 & 17,558\\
Hợp mọi peak & 0,6625 & 76,140\\
\bottomrule
\end{tabular}
\end{table}
Bootstrap ghép cặp cho chênh F1 cross-attention trừ SWT là $-0{,}0276$, khoảng 95\% $[-0{,}0437;-0{,}0144]$. Cặp FastQRS+SWT trừ SWT là $-0{,}0008$, khoảng $[-0{,}0064;+0{,}0045]$. Phép thử chưa chứng minh fusion cải thiện detector đơn; do đó em không đưa fusion vào pipeline chính. Giới hạn exposure của checkpoint SWT và quy mô crop vẫn phải được giữ khi diễn giải.
\artifact{experiment/report/benchmark_runs/qrs_detector_fusion_crossattn_20261027/REPORT.md}

\section{Thử nghiệm suy luận Construe}
Nhánh host Construe được thử trên 10 đoạn AFDB và 6 đoạn MITDB dài 61,44 giây, chấm phần giữa để tránh biên. QRS F1 của production sau Construe giảm từ 0,9913 xuống 0,9481 trên nhóm AFDB nhưng tăng từ 0,8443 lên 0,8914 trên MITDB. Phần AF chỉ có 10 quyết định, quá nhỏ để suy ra hiệu năng toàn AFDB. Chi phí suy luận và trường hợp cạn search làm phương pháp này phù hợp hơn với baseline nghiên cứu/phân tích lỗi trên host ở trạng thái hiện tại.
\artifact{experiment/src/host_bench/tomas_construe/README.md}

\chapter{Kết quả đánh giá chất lượng tín hiệu}
\section{Benchmark BUTQDB}
BUTQDB cung cấp nhãn chất lượng theo thời gian. Em kiểm kê 18 record của 15 người với tổng thời lượng ECG 478,73 giờ; phần có đủ consensus để chấm gồm 35.773 cửa sổ 10 giây, tương đương 99,37 giờ. Lớp 1 và 2 được gộp là đủ cho QRS; lớp 3 là xấu. Nhãn cửa sổ lấy lớp xấu nhất trong khoảng; phần chưa được gán nhãn bị loại theo protocol. Đánh giá dùng 5 fold theo người, với tiền xử lý và huấn luyện chỉ trên phần train.

\begin{table}[H]
\centering\small
\caption{BUTQDB: metric SQI, không phải metric AF}
\begin{tabular}{lrrrr}
\toprule Cấu hình & BA & AUC & Bao phủ & Xấu trong phần nhận\\\midrule
Orphanidou adaptation & 0,9249 & 0,9351 & 71,86\% & 0,12\%\\
Ho[NeuroKit] ít nhất 80\% & 0,9179 & 0,9792 & 72,79\% & 0,58\%\\
Kramer[ECGAssess] & 0,8935 & 0,8936 & 74,52\% & 1,82\%\\
Zhao[NeuroKit] & 0,5051 & 0,8020 & 88,82\% & 15,90\%\\
Project-5SQI Logistic & 0,9647 & 0,9941 & 83,33\% & 1,02\%\\
bSQI rồi Zhao cascade & 0,5518 & 0,5518 & 8,70\% & 0,00\%\\
\bottomrule
\end{tabular}
\end{table}
Project-5SQI có điểm gộp cao, nhưng ba record chiếm 88,58\% số cửa sổ. Khoảng bootstrap theo người của chênh BA so với Orphanidou/Ho chứa không; em chưa kết luận mô hình tốt hơn ổn định ở mức người. Cascade có tỷ lệ xấu bằng không trong phần nhận nhưng chỉ giữ 8,70\%, nên tỷ lệ này cần đọc cùng coverage.

Trong lượt đánh giá này, BUTQDB không có nhãn R-peak/AF độc lập tương ứng để đo BPM hoặc AF. Không thể dùng nhãn chất lượng thay cho nhãn nhịp bệnh lý.
\artifact{experiment/report/benchmark_runs/butqdb_20261020/REPORT.md}

\section{Các thử nghiệm quality gate chưa đạt}
Ở nhánh QTDB+MA trước đó, em đánh giá 17.220 cửa sổ, gồm 10.500 development và 6.720 test theo record nguồn. Nhãn tốt được xây từ yêu cầu không FP/FN QRS, có đủ cặp RR liên tiếp khớp và RR MAE không quá 20 ms. SVM fusion đạt AUC 0,8554, BA 0,8008, nhưng ngưỡng kiểm soát selective risk chỉ giữ 0/6.720 cửa sổ test. Do đó không có gate hữu dụng để export từ kết quả này.

\begin{table}[H]
\centering\small
\caption{Một số gate học máy lịch sử và mục tiêu risk không quá 5\%}
\begin{tabular}{lrrr}
\toprule Mô hình & Test AUC & Bao phủ & Xấu trong phần nhận\\\midrule
Compact Extra Trees, 5 SQI & 0,818 & 22,40\% & 17,87\%\\
Compact Logistic & 0,801 & 4,99\% & 5,97\%\\
Full Extra Trees & 0,812 & 17,31\% & 19,69\%\\
vital\_sqi statistical & 0,772 & 10,52\% & 23,34\%\\
FW-5 morphology & 0,780 & 5,43\% & 5,48\%\\
\bottomrule
\end{tabular}
\end{table}
Các mô hình đều chưa đạt mục tiêu đã đặt hoặc có coverage quá thấp. FW-5 còn có giới hạn feature cache chưa được chứng minh bit-exact với streaming. Em giữ các kết quả âm này trong báo cáo vì chúng là căn cứ không đưa một gate học máy chưa đạt xuống firmware.
\artifact{experiment/report/outputs/sqi_fusion_protocol_v1/REPORT.md}
\artifact{experiment/report/outputs/sqi_firmware_native_v1/REPORT.md}

\section{CinC2011 và phạm vi parity}
CinC2011 có nhãn mức chấp nhận của toàn bản ghi 12 lead, không có nhãn chất lượng riêng cho Lead II. Vì vậy, phép dự đoán nhãn record từ một lead không thể được dùng như bảng xếp hạng SQI một chuyển đạo. Artifact accuracy lịch sử cũng không còn đầy đủ; em ưu tiên BUTQDB cho kết luận chất lượng.

Tuy vậy, dùng cùng waveform Lead II và cùng peak để so Python/C vẫn hữu ích cho kiểm tra chuyển mã. Zhao-simple đạt đồng thuận score/category trên 998 cặp, Kramer có 5 khác biệt, còn Clifford được kiểm tra sai số năm feature. Với Orphanidou, con số đồng thuận 99,90\% chỉ thuộc kernel RR-only; không được gán cho toàn runtime có template. Ho chỉ có kernel ICI được port; hai detector UNSW/Khamis và NeuroKit chưa được port cùng.
\artifact{experiment/report/benchmark_runs/sqi_c_port_parity_cinc2011_20261022/REPORT.md}
\artifact{experiment/report/CLIFFORD_RAHMAN_SQI_IMPLEMENTATION_REPORT.md}

\section{SQI trên dữ liệu thiết bị}
Các CSV thiết bị đã được xử lý bằng NeuroKit tham chiếu. Một cohort 22 CSV có 79/122 khối 10 giây được nhận (64,8\%). Cohort AD8232 47 CSV sau regrid 360 Hz có 260/267 khối được nhận (97,4\%), còn bản rolling 10 giây/hop 1 giây có 2.259/2.318 khối được nhận (97,5\%). Những tỷ lệ này là đầu ra của thuật toán, không phải accuracy do chưa có nhãn chất lượng độc lập.

Không thể diễn giải 64,8\% tăng thành 97,4\% là một cải thiện phương pháp: cohort và tiền xử lý đã khác. Nội suy timestamp giúp đưa dữ liệu lên lưới chung nhưng không sửa clipping; phép Ho toàn record/per-RR là phân tích hồi cứu, không tự chứng minh hành vi nhân quả trên ESP32.
\artifact{experiment/report/benchmark_runs/sqi_ho_neurokit_ad8232_rolling_20261027/REPORT.md}

\chapter{Kết quả phát hiện AF và lựa chọn pipeline}\label{chap:af}
\section{Mô hình RR và ảnh hưởng của lượng tử hóa}
Benchmark mô hình ban đầu có 78.921 cửa sổ của 25 record, split 18/7 theo record, trong đó test có 20.194 cửa sổ. Năm feature được chọn chỉ trên train bằng 5-fold StratifiedGroupKFold: lấy tập nhỏ nhất có BA cross-validation nằm trong 0,5 điểm phần trăm so với kết quả tốt nhất. Với đặc trưng rút gọn từ annotation RR, Extra Trees float đạt BA 0,9158 và F1 0,8942; Logistic Regression đạt BA 0,9022; Random Forest 0,8664; XGBoost 0,8244. Đây là metric theo cửa sổ, không phải metric giây của pipeline liên tục.
\artifact{experiment/report/outputs/af_detection/final_benchmark/REPORT.md}

\begin{table}[H]
\centering\small
\caption{So sánh float/INT8 trên 20.194 cửa sổ test phía máy chủ}
\begin{tabular}{lrrrrr}
\toprule Mô hình & BA float & BA INT8 & AUC float & AUC INT8 & Đồng thuận\\\midrule
Logistic Regression & 0,90221 & 0,89822 & 0,98315 & 0,97749 & 98,98\%\\
Random Forest & 0,86641 & 0,86466 & 0,98092 & 0,98013 & 99,47\%\\
Extra Trees & 0,91576 & 0,91291 & 0,98166 & 0,97908 & 99,47\%\\
\bottomrule
\end{tabular}
\end{table}
Lượng tử hóa gây suy giảm nhỏ trong benchmark này. Host C được kiểm tra parity với bản lượng tử hóa Python trên toàn test. Các số liệu này không được chuyển thành tuyên bố accuracy đo trên board chỉ vì cùng mô hình đã chạy kernel trên ESP32.

\section{AFDB waveform timeline với các mô hình đóng băng}
Em phát lại sáu waveform test, 60,393 giờ, trên tín hiệu gốc và 18 điều kiện nhiễu. Nhãn được lấy tại thời điểm lưới 1 giây; các giây chuyển nhịp không bị bỏ. Đầu ra chỉ có hiệu lực khi đã sẵn có theo processing time.

\begin{table}[H]
\centering\small
\caption{AFDB gốc, đầu ra thô, không SQI gate; metric theo giây}
\begin{tabular}{lrrrr}
\toprule Mô hình & AUC & BA & Bao phủ & Recall AF hiệu dụng\\\midrule
Logistic Regression & 0,9767 & 0,8883 & 98,59\% & 80,44\%\\
Random Forest & 0,9764 & 0,8478 & 98,59\% & 72,27\%\\
Extra Trees float & 0,9729 & 0,8926 & 98,59\% & 82,96\%\\
XGBoost & 0,9737 & 0,8234 & 98,59\% & 68,19\%\\
Extra Trees INT8 C & 0,9708 & 0,8899 & 98,59\% & 82,82\%\\
\bottomrule
\end{tabular}
\end{table}
Các mô hình được huấn luyện trước từ annotation RR; lượt waveform dùng RR do detector production đo. Như vậy, sai số QRS và thay đổi lịch phát cửa sổ trở thành một phần của hiệu năng AF. Không thể giữ số accuracy của annotation RR rồi gọi đó là độ chính xác hệ thống.

\section{Đánh giá trên LTAFDB và SHDB-AF}
Đánh giá ngoài gồm 182 waveform đã kiểm tra checksum, không refit, không chọn lại feature hoặc ngưỡng. AF/AFIB là dương tính; các rhythm khác có annotation rõ, gồm AFL hoặc AB, là âm tính. Thời gian trước annotation rhythm đầu tiên là unknown.

\begin{table}[H]
\centering\small
\caption{External waveform, đầu ra thô không gate}
\begin{tabular}{llrrr}
\toprule Dataset & Mô hình & BA & AUC & Recall AF hiệu dụng\\\midrule
LTAFDB & Logistic Regression & 0,9184 & 0,9504 & 90,55\%\\
LTAFDB & Random Forest & 0,9258 & 0,9663 & 93,05\%\\
LTAFDB & Extra Trees float & 0,9193 & 0,9567 & 95,04\%\\
LTAFDB & XGBoost & 0,9198 & 0,9651 & 92,52\%\\
LTAFDB & Extra Trees INT8 & 0,9120 & 0,9545 & 95,24\%\\\midrule
SHDB-AF & Logistic Regression & 0,8956 & 0,9228 & 86,66\%\\
SHDB-AF & Random Forest & 0,8953 & 0,9481 & 88,65\%\\
SHDB-AF & Extra Trees float & 0,8929 & 0,9418 & 89,97\%\\
SHDB-AF & XGBoost & 0,8927 & 0,9485 & 88,51\%\\
SHDB-AF & Extra Trees INT8 & 0,8871 & 0,9401 & 90,18\%\\
\bottomrule
\end{tabular}
\end{table}
Coverage lần lượt là 97,52\% và 99,18\%. Kết quả cho thấy xếp hạng BA và recall không hoàn toàn giống nhau; chọn mô hình theo một metric duy nhất có thể thay đổi đánh đổi giữa bỏ sót và cảnh báo sai.

\section{Gate SQI có thể loại cả AF thật}
Trên AFDB gốc với Extra Trees float, Orphanidou gate làm BA tăng từ 0,8926 lên 0,9122, nhưng coverage giảm từ 98,59\% xuống 85,75\% và tỷ lệ giữ lại thời gian AF chỉ 74,03\%. Ho[NeuroKit] ít nhất 80\% giữ 86,05\% thời gian AF, BA 0,8930 và coverage 86,86\%.

Trên LTAFDB, Orphanidou gate làm BA Extra Trees float tăng 0,9193 lên 0,9482, nhưng recall episode ở đầu ra thô giảm 80,38\% xuống 61,69\%. Đây là lý do em không chọn gate cứng chỉ dựa trên BA sau khi từ chối các đoạn khó. Gate ở các thử nghiệm này chỉ điều khiển việc chấp nhận đầu ra; trạng thái AF vẫn tiếp tục chạy và không có phép đo tiết kiệm năng lượng đi kèm.

\section{So sánh annotation RR với FastQRS}
Thử nghiệm đầu tiên giữ cùng mô hình, nhãn và độ trễ gom khối 10 giây. Với INT8 thô, nhánh annotation đạt BA 0,9105, F1 0,8787 và cảnh báo sai 3,70/giờ non-AF; FastQRS đạt BA 0,8877, F1 0,8490 và 4,98/giờ. Recall AF hiệu dụng lần lượt 83,60\% và 83,30\%.

Đối chứng annotation còn có các khoảng thiếu mốc nhịp làm runtime giữ đầu ra cũ, tuổi tối đa khoảng 835,908 giây. Khi loại đầu ra cũ trên 30 giây trong phân tích độ nhạy, xu hướng BA thấp hơn của FastQRS vẫn còn. Mốc .qrs là tham chiếu tự động chưa hiệu đính toàn bộ, không phải cận trên lý tưởng đảm bảo đúng mọi nhịp.

\section{Lựa chọn detector bằng quy tắc đã khóa}
Để đi từ so sánh đơn lẻ đến quyết định cấu hình, em mở rộng thành năm detector thực: production, FastQRS, WIM Nguyen, WIM Elgendi và XQRS C 250 Hz. Mọi ứng viên dùng cùng notch đầu vào và cùng Extra Trees INT8. Cấu hình production trì hoãn 10 giây được thêm như một đối chứng độ trễ, không phải detector thứ sáu.

Sàng lọc gồm 114 job AFDB và 84 record LTAFDB. Quy tắc theo đầu ra xác nhận 3/4 yêu cầu recall AF hiệu dụng không giảm quá 2 điểm phần trăm, cảnh báo sai không tăng quá 0,5/giờ và coverage không giảm quá 2 điểm phần trăm. Trong số ứng viên đạt các điều kiện này, chọn BA cao nhất. Nếu không có ứng viên đạt, chọn BA cao nhất làm challenger để đối chiếu, nhưng chưa coi là vượt sàng lọc.

\begin{table}[H]
\centering\small
\caption{LTAFDB: sàng lọc detector, đầu ra xác nhận 3/4}
\begin{tabular}{lrrrr}
\toprule Detector & BA & \shortstack{Recall AF\\hiệu dụng} & Bao phủ & \shortstack{Cảnh báo sai bắt đầu\\/giờ non-AF}\\\midrule
Production liên tục & 0,9233 & 94,12\% & 97,52\% & 2,7953\\
FastQRS & 0,9025 & 93,27\% & 98,22\% & 4,3395\\
WIM Nguyen & 0,9212 & 91,08\% & 95,43\% & 2,9269\\
WIM Elgendi & 0,9037 & 85,10\% & 91,33\% & 2,6459\\
XQRS C 250 Hz & 0,9103 & 87,12\% & 91,64\% & 2,7543\\
\bottomrule
\end{tabular}
\end{table}
Không ứng viên nào đạt đầy đủ điều kiện sàng lọc. WIM Nguyen có BA cao nhất trong các ứng viên nên được khóa làm challenger theo quy tắc dự phòng. Nó giảm recall hiệu dụng khoảng 3,05 điểm phần trăm và coverage 2,09 điểm trên LTAFDB, do đó không được mô tả là đã vượt sàng lọc.

\begin{table}[H]
\centering\small
\caption{Pipeline production + Extra Trees INT8, xác nhận 3/4}
\begin{tabular}{lrrrrr}
\toprule Dataset & BA & F1 & Bao phủ & \shortstack{Recall AF\\hiệu dụng} & \shortstack{Cảnh báo sai bắt đầu\\/giờ non-AF}\\\midrule
AFDB gốc & 0,8785 & 0,8449 & 98,59\% & 79,44\% & 1,9226\\
LTAFDB & 0,9233 & 0,9329 & 97,52\% & 94,12\% & 2,7953\\
SHDB-AF & 0,8898 & 0,7596 & 99,18\% & 88,82\% & 2,2224\\
\bottomrule
\end{tabular}
\end{table}
Trên SHDB-AF, Nguyen đạt BA 0,8889, recall AF hiệu dụng 87,90\% và 2,1445 cảnh báo sai/giờ non-AF. Bootstrap ghép cặp 5.000 lần theo 93 bệnh nhân cho các chênh lệch ở Bảng~\ref{tab:confirm}. Dấu chênh lệch được tính bằng Nguyen trừ production.

Quy tắc chấp nhận thay thế còn yêu cầu ứng viên đã vượt sàng lọc, cận dưới khoảng bootstrap 95\% của $\Delta BA>0$, cận dưới của $\Delta Recall_{\mathrm{AF,eff}}\geq-0{,}02$, cận trên của $\Delta$ số cảnh báo sai/giờ non-AF $\leq+0{,}5$ và mức giảm coverage tại ước lượng điểm không quá $0{,}02$. Trên AFDB có nhiễu, mức giảm recall hiệu dụng trung bình không quá $0{,}02$ và không điều kiện nhiễu riêng nào giảm quá $0{,}05$. Đây là các biên quyết định kỹ thuật được khai báo trước cho lượt thử, chưa phải ngưỡng chấp nhận lâm sàng.

\begin{table}[H]
\centering\small
\caption{Xác nhận SHDB-AF theo bệnh nhân, cùng ngưỡng đã khóa}\label{tab:confirm}
\begin{tabular}{lrr}
\toprule Chỉ số & Chênh lệch & Khoảng bootstrap 95\%\\\midrule
BA & $-0{,}000947$ & $[-0{,}002360;+0{,}000548]$\\
Recall AF hiệu dụng & $-0{,}009215$ & $[-0{,}019664;-0{,}003118]$\\
Cảnh báo sai/giờ non-AF & $-0{,}077886$ & $[-0{,}356029;+0{,}093642]$\\
Coverage & $+0{,}002845$ & $[-0{,}002484;+0{,}011413]$\\
\bottomrule
\end{tabular}
\end{table}
BA chưa được chứng minh cao hơn production. Mặt khác, ở MA $-6$ dB trên AFDB, recall AF hiệu dụng của production là 85,11\%, còn Nguyen 67,25\%, giảm 17,86 điểm phần trăm. Điều kiện độ bền nhiễu được khai báo trước cũng chưa đạt.

Kết luận lựa chọn là \textbf{giữ detector firmware liên tục kết hợp Extra Trees INT8 và xác nhận 3/4}. Đây là quyết định trong shortlist đã thử; WIM Zhai và các mô hình QRS học sâu chưa nằm trong đợt detector-selection này. Em không suy rộng thành chứng minh tối ưu trên mọi thuật toán.

Audit độc lập đã kiểm tra 296 job, 1.498 replay C mới, 112.826.672 lượt đỉnh R, 176.090.872 ô timeline, 72.128 kiểm tra counts theo record và 8.544 kiểm tra counts gộp. Hash source/input được đối chiếu, thời điểm khóa lựa chọn có trước confirmation và bootstrap được tái lập. Tất cả đều đạt trong phạm vi artifact đã kiểm tra.
\artifact{experiment/report/benchmark_runs/af_detector_selection_20261031/decision.json}
\artifact{experiment/report/benchmark_runs/af_detector_selection_20261031/independent_audit.json}

\section{Thử nghiệm mô hình waveform AF}
Hai thử nghiệm waveform sử dụng 7.499 cửa sổ từ 177 nhóm bệnh nhân LTAFDB+SHDB-AF: 5.895 train và 1.604 validation, tách nhóm 141/36. Sáu waveform AFDB cung cấp 21.616 cửa sổ 10 giây thuần nhãn để thử chuyển miền. Chỉ cửa sổ thuần được chấm, nên bảng này không được xếp hạng trực tiếp với metric timeline RR.

\begin{table}[H]
\centering\small
\caption{AFDB cửa sổ waveform thuần nhãn, ngưỡng 0,50}
\begin{tabular}{lrrrrr}
\toprule Mô hình & BA & Độ nhạy & Đặc hiệu & F1 & AUC\\\midrule
ECGFounder đóng băng + logistic & 0,9035 & 0,8621 & 0,9450 & 0,8666 & 0,9492\\
ConvNeXt FP32 & 0,7869 & 0,7283 & 0,8455 & 0,6982 & 0,8904\\
ConvNeXt TFLite lượng tử hóa & 0,7616 & 0,6774 & 0,8459 & 0,6660 & 0,8747\\
\bottomrule
\end{tabular}
\end{table}
Encoder FP32 ECGFounder trích xuất có kích thước 123.048.857 B; đây chưa phải mô hình đã lượng tử hóa cho điện thoại. ConvNeXt có 766.818 tham số, artifact TFLite 1.125.000 B, đầu vào/đầu ra INT8 nhưng vẫn còn 41 phép tính float. Độ trễ điện thoại và tích hợp ứng dụng chưa được đo. Kết quả ECGFounder gộp cần đọc cùng per-record: có record ít cửa sổ AF và F1 rất thấp, nên chưa kết luận tính ổn định từ BA gộp.

\chapter{Giải thích mô hình và phân tích lỗi}
\section{Thiết kế tập giải thích}
Em dùng exact TreeSHAP để phân tích 50 cửa sổ được chọn có chủ đích: 15 TP, 15 TN, 10 FP và 10 FN từ các waveform test có sẵn. Mô hình, bộ feature và ngưỡng không được fit lại khi tạo giải thích. Tập này phục vụ phân tích lỗi, không phải mẫu ngẫu nhiên để ước lượng lại accuracy.

Sai số additivity lớn nhất là $2{,}25\times10^{-13}$, kiểm tra rằng tổng đóng góp cùng giá trị nền khớp đầu ra của mô hình đang giải thích. Đây là kiểm tra tính nhất quán toán học của explanation, không phải kiểm định ý nghĩa sinh lý.

\begin{table}[H]
\centering
\caption{Độ lớn đóng góp SHAP trung bình trên 50 mẫu chọn chủ động}
\begin{tabular}{lr}
\toprule Đặc trưng & Trung bình $|SHAP|$ xấp xỉ\\\midrule
pNN50 & 0,227\\
$SD_2$ & 0,089\\
Median RR & 0,070\\
Shannon entropy & 0,063\\
$SD_1/SD_2$ & 0,016\\
\bottomrule
\end{tabular}
\end{table}
Các FP thường có đóng góp dương mạnh từ pNN50 và $SD_2$: RR không đều ở nhịp không AF có thể bị mô hình gán AF. Một số FN có pNN50/entropy kéo điểm xuống. Phân tích cũng cho thấy lỗi tập trung ở một vài record, củng cố yêu cầu báo cáo theo record/bệnh nhân thay vì chỉ theo tổng cửa sổ.

\section{Giới hạn fidelity và đầu ra giải thích}
TreeSHAP đang giải thích Extra Trees float. Mô hình INT8 có tỷ lệ đồng thuận cao trên tập test đã đo nhưng không giống float ở mọi đầu vào. Do đó, nếu server giải thích bản float cho một kết quả device INT8, cần lưu cả hai điểm và đánh dấu bất đồng; không gọi explanation float là giải thích chính xác quyết định INT8 trong trường hợp đó.

Một đầu ra mong muốn gồm waveform có đánh dấu R, RR, năm feature, điểm AF, trạng thái xác nhận, SQI, timestamp và phiên bản mô hình. Phần văn bản diễn giải cần gắn vào bằng chứng đó. Hiện em đã có phân tích ngoại tuyến và các hình TP/TN/FP/FN; phần dịch vụ tự động tạo báo cáo cho dữ liệu mới chưa được xác nhận hoạt động end-to-end.
\artifact{experiment/report/AF_XAI_50_REPORT.md}

\chapter{Khả năng triển khai và số đo ESP32}
\section{Detector QRS trên board}
Các benchmark thực được đo trên ESP32-D0WD-V3 ở 160 MHz. Đã có 15 hiện thực C hoàn tất cohort 59 đoạn, tương ứng 885 kết quả phương pháp--cửa sổ. Bảng~\ref{tab:deviceqrs} trình bày thời gian thực thi detector và workspace của từng phương pháp, chưa phải độ trễ từ thu nhận ECG tới hiển thị kết quả.

\begin{table}[H]
\centering\small
\caption{MITDB crop trên ESP32: thời gian xử lý một khối 10 giây}\label{tab:deviceqrs}
\begin{tabular}{lrrr}
\toprule Hiện thực C & F1 $\pm50$ ms & p95 (ms) & Workspace (B)\\\midrule
Production & 0,9371 & 29,639 & 0 bổ sung\\
WIM Nguyen & 0,9846 & 9,976 & 68.656\\
FastQRS & 0,9914 & 25,350 & 65.196\\
WIM Elgendi & 0,9558 & 13,949 & 61.588\\
WIM Zhai & 0,9889 & 180,455 & 58.744\\
XQRS C 250 Hz & 0,9906 & 40,349 & 59.248\\
GQRS C & 0,9218 & 89,174 & 67.288\\
Pan--Tompkins++ & 0,9789 & 278,630 & 77.576\\
WIM Xia & 0,8038 & 1.538,483 & 73.192\\
\bottomrule
\end{tabular}
\end{table}
Con số workspace bổ sung bằng không của production không có nghĩa toàn thuật toán không dùng RAM; nó phản ánh cách benchmark hạch toán buffer ngoài. XQRS 250 Hz chưa cộng bước resample vào timing. Xia hoàn tất nhưng có watchdog warning, còn Xu chỉ trả được tám cửa sổ trước khi dừng; không có full-cohort metric Xu để báo cáo. Kích thước toàn image benchmark không phải Flash tăng thêm của từng phương pháp.

\section{Classifier AF INT8 trên board}
\begin{table}[H]
\centering\small
\caption{Kernel phân loại AF: vector INT8 cố định, 10.000 lần gọi}
\begin{tabular}{lrrrr}
\toprule Mô hình & p50 ($\mu$s) & Max ($\mu$s) & Model Flash (B) & Working RAM (B)\\\midrule
Logistic Regression & 8,160 & 9,730 & 43 & 17\\
Random Forest & 607,320 & 620,620 & 215.833 & 13\\
Extra Trees & 391,580 & 405,550 & 141.881 & 13\\
\bottomrule
\end{tabular}
\end{table}
Các phép đo gồm 101 batch timing và không bao gồm trích xuất feature. Working RAM ở bảng là các giá trị làm việc của API inference, không phải toàn bộ RAM AF/runtime/pipeline. Image benchmark Extra Trees đo lại là 294.416 B; số 157.424 B trong profile cũ thuộc image khác và đã được thay thế. Không suy ra công suất hoặc accuracy trên thiết bị từ thời gian kernel.
\artifact{experiment/report/benchmark_runs/device_bench_20261023_sqi_af/REPORT.md}

\section{SQI C port: phạm vi và chi phí}
\begin{table}[H]
\centering\small
\caption{Phạm vi SQI và số đo thực trên ESP32}
\begin{tabularx}{\textwidth}{L{3.2cm}L{3cm}>{\raggedright\arraybackslash}X}
\toprule Khối & Timing được đo & Phạm vi phải giữ khi diễn giải\\\midrule
Orphanidou rules + template & Finalize 18,635 ms, một lần gọi & State 7.372 B và template 4.320 B; detector đầu vào không thuộc phép đo này.\\
Kramer rules-only & Push p99 14,031 $\mu$s/mẫu & State 408 B; khác bản waveform đầy đủ.\\
Kramer waveform full & Finalize 19,184 s, một lần gọi & Có ba watchdog warning, vẫn trả kết quả; chưa phù hợp nhịp gọi 10 giây trong hiện thực này.\\
Zhao waveform scorer & Finalize 1,830 s, một lần gọi & State 14.532 B, scratch benchmark riêng; cần xét chi phí khi tích hợp.\\
Ho ICI kernel & p50 51,806 $\mu$s & Chỉ kernel cleanup/matching/decision; thiếu hai QRS detector của toàn phương pháp.\\
Clifford 5 feature & p50 119,612 ms & Workspace FFT 65.536 B; chưa gồm đầy đủ bSQI/iSQI/data fusion.\\
\bottomrule
\end{tabularx}
\end{table}
Các finalize đo một lần không được gọi là percentile. Các số đo kernel có input tổng hợp hoặc cố định, không phải độ chính xác SQI trên ECG người. Sự khác nhau giữa một port feature, một kernel quyết định và toàn phương pháp là phần quan trọng của đánh giá khả năng triển khai.

\section{Mô hình học sâu và giới hạn tensor}
TFLite Micro đã được đưa vào cấu hình ESP-IDF và ví dụ hello\_world lượng tử hóa chạy được trên board. Đây là kiểm tra toolchain/runtime, chưa phải mô hình ECG. Với các mô hình ECG, kết quả chuyển đổi và chạy thực tế khác nhau:
\begin{itemize}
\item Han CNN chuyển INT8 thành công với artifact 978.888 B, nhưng image thử không cấp phát được một arena liên tục 266.240 B; inference chưa chạy.
\item Laitala full-INT8 gặp lỗi calibrator native dưới TensorFlow 2.21/Windows. Bản dynamic-range có activation lớn nhất 512.000 B nên chưa phù hợp bộ nhớ trong đang có.
\item ECG2RR dynamic-range cũng có activation lớn 512.000 B; phép chuyển đổi không tự tạo ra metric QRS trên board.
\item Han RNN full-INT8 không tạo được model; SWT--SepConv chưa có TFLM export/port được xác nhận.
\end{itemize}
Những kết quả này thuộc các định dạng và cấu hình đã thử, không phải chứng minh mọi cách lượng tử hóa hoặc thiết kế lại đều bất khả thi. Kích thước file trọng số hoặc dung lượng toàn chip Flash chưa đủ để kết luận mô hình chạy được; activation, operator, arena liên tục và partition cũng cần kiểm tra.

\section{Kiểm tra đường dữ liệu và giới hạn phép đo}
Các smoke SD/BLE lịch sử xác nhận ghi segment, quay vòng file, sequence và event. Một số đợt UART được thực hiện khi frontend chưa kết nối, nên BPM/AF trong những capture đó không có ý nghĩa accuracy trên người. Dữ liệu AD8232 về sau đã tồn tại, nhưng chưa có bộ nhãn độc lập tương ứng để xác nhận sàng lọc AF trên người đeo.

Em chưa có số đo mới của toàn pipeline liên tục ứng với đợt detector-selection ngày 31/10. Không cộng latency các benchmark riêng để gọi là latency end-to-end, và không xem việc image production từng được khôi phục/boot ở lượt trước như bằng chứng trạng thái board hiện giờ.

\chapter{Thảo luận, giới hạn và hướng phát triển}
\section{Những kết quả đã đạt được}
Đề tài đã hình thành một chuỗi thực nghiệm có thể truy vết từ dữ liệu tới quyết định cấu hình. Em đã khảo sát các nhóm QRS cổ điển, biến đổi và học sâu; tổ chức phép common-input và toàn record; chuyển nhiều thuật toán sang C; đo kernel trên ESP32; đánh giá SQI trên nhãn chất lượng; triển khai AF INT8 trong source; thử mô hình waveform; và phân tích lỗi bằng TreeSHAP.

Giá trị chính của kết quả không nằm ở việc một thuật toán có số F1 cao nhất trên một bảng. Các thí nghiệm cho thấy ba vấn đề: sai số peak có thể lan truyền sang AF; gate có thể tăng BA nhưng giảm lượng AF được phục vụ; và một mô hình có file nhỏ hoặc đã chuyển đổi chưa chắc chạy được trên MCU. Việc giữ production sau khi một số ứng viên chưa đạt cũng là một quyết định có bằng chứng.

\section{Giới hạn của bằng chứng}
\begin{enumerate}
\item AFDB đã xuất hiện trong phát triển; nhiều waveform train không có sẵn để thử một chương trình huấn luyện lại công bằng. Đánh giá hiện tại dùng mô hình đóng băng.
\item Annotation .qrs của AFDB là tự động; nó hỗ trợ đối chiếu nhưng không phải đỉnh R được chuyên gia duyệt toàn bộ. Nhãn rhythm .atr và nhãn peak cần được phân biệt.
\item Cohort QTDB MLII ở common-input chồng nguồn với MITDB. BUTQDB bị lệch số cửa sổ giữa record; số đo gộp không thay cho kiểm tra theo người.
\item Các mô hình waveform chấm cửa sổ thuần và loại chuyển nhịp; metric của chúng chưa thay cho metric timeline của detector RR.
\item Một số kiểm tra phần cứng là kernel hoặc smoke ngắn; chưa có endurance, năng lượng và accuracy thực địa cho toàn thiết kế hai tầng.
\item Chưa có bằng chứng MECG-E, fusion AF và XAI service đã tích hợp. Chưa có thử nghiệm tiến cứu chứng minh hiệu quả hoặc mức sẵn sàng lâm sàng.
\end{enumerate}

\section{Hướng phát triển có thể kiểm chứng}
Ưu tiên tiếp theo là bổ sung bộ dữ liệu thiết bị có nhãn R-peak/chất lượng/AF độc lập và provenance rõ, sau đó đo pipeline tích hợp về latency, dropped samples, stack, heap và năng lượng. Khi cân nhắc TTL, cần khóa policy trước và đánh giá lại recall/coverage trên toàn thời gian; không chọn timeout chỉ vì một bảng sau lọc đẹp hơn.

Đối với SQI, em đề xuất đánh giá chính sách điều phối thay vì chỉ một gate cứng: kết quả đủ chất lượng được xuất tại chỗ, trường hợp còn khả năng phân tích nhưng bất định được gửi waveform để kiểm tra bổ sung, và trường hợp không dùng được yêu cầu đo lại. Đây là giả thuyết thiết kế, cần đo tỷ lệ gọi server, cảnh báo sai, recall và độ trễ; chưa phải kết quả đã đạt.

Với server, cần khóa preprocessing, version model và cách đối chiếu kết quả float/INT8. Nếu thử fusion, phải có detector đơn/mô hình đơn mạnh làm đối chứng và kiểm tra trên tập ngoài được giữ riêng. Với khử nhiễu, cần mở rộng CleanGuard từ smoke sang cohort đầy đủ, xác định annotation phù hợp và đánh giá ảnh hưởng lên AF, trước khi nối MECG-E vào đường chính.

\chapter{Kết luận}
Trong phạm vi khóa luận, em đã xây dựng và kiểm tra nhiều thành phần của hệ thống ECG nhúng, đồng thời khảo sát các phương pháp làm nền cho hướng mở rộng. Các kết quả hiện có hỗ trợ việc giữ pipeline detector firmware liên tục, tính năm feature từ 30 RR, chạy Extra Trees INT8 và xác nhận 3/4. Cấu hình này được lựa chọn từ phép so sánh downstream có ma trận nhiễu và bước đối chiếu theo bệnh nhân, thay vì chỉ dựa vào độ chính xác QRS hoặc BPM.

SQI, mô hình waveform và XAI đã có những kết quả độc lập hữu ích, nhưng mức tích hợp và phạm vi kiểm chứng khác nhau. Em giữ rõ sự khác biệt giữa phần đã chạy, phần chưa đạt và phần dự kiến, để các bước phát triển tiếp theo có thể kiểm tra bằng cùng nguyên tắc về dữ liệu, thời gian sẵn có, coverage và bằng chứng triển khai.

\appendix
\chapter{Bản đồ mã nguồn và artifact}
\section{Các báo cáo đọc chính}
\begin{itemize}
\item QRS/RR/BPM: \path{experiment/report/qrs_rr_bpm/REPORT.md}.
\item AF, model waveform và lựa chọn detector: \path{experiment/report/af_det/report.md}.
\item Tổng hợp metric: \path{experiment/report/REPORT_METRICS.md}.
\item Phạm vi port: \path{experiment/report/benchmark_runs/C_PORT_DEVICE_STATUS.md}.
\item Sơ đồ hệ thống: \path{experiment/report/af_det/af_system_architecture.drawio}.
\end{itemize}

\section{Source và dữ liệu thực nghiệm}
\begin{longtable}{L{3.6cm}L{11cm}}
\toprule Hạng mục & Vị trí tương đối trong workspace\\\midrule\endfirsthead
\toprule Hạng mục & Vị trí tương đối trong workspace\\\midrule\endhead
Firmware & \path{firmware/esp32/main/}\\
QRS production & \path{firmware/esp32/main/task/Core1/ECGPipeline/ExFeature/ExtractFeature.c}\\
AF runtime/model & \path{firmware/esp32/main/components/af_model/}\\
SQI firmware & \path{firmware/esp32/main/components/sqi_gate/}\\
CleanGuard runner & \path{experiment/src/experiments/qtdb_cleanguard_rpeaks.py}\\
Metric/baseline khử nhiễu & \path{experiment/src/ecg_benchmark/}\\
AF feature, training, XAI & \path{experiment/src/experiments/af_detection/}\\
Runner lịch sử & \path{experiment/src/esp32_archive/experiment_code_and_assets/tools/}\\
Repository khảo sát & \path{external/} và các thư mục external được lưu trong archive.\\
Dữ liệu cục bộ & \path{experiment/data/}; các dữ liệu lớn được bỏ qua trong Git.\\
\bottomrule
\end{longtable}

\section{Các lượt chạy được dùng trong báo cáo}
Các đường dẫn dưới đây đều tương đối với \path{experiment/report/benchmark_runs/}:
\begin{itemize}
\item \path{wim_common_input_mitdb_qtdb_20261024/} và \path{qrs_common_input_mitdb_qtdb_20261024_rerun/}: QRS common-input, vector, UART và số đo board.
\item \path{qrs_detector_fusion_crossattn_20261027/}: folds, OOF, bootstrap và metric fusion.
\item \path{butqdb_20261020/}: SQI theo subject và các cấu hình tham chiếu.
\item \path{af_timeline_20261021/}, \path{af_external_20261022/}: AF đóng băng và output-gate.
\item \path{af_fastqrs_e2e_20261031/}: annotation RR so với FastQRS và audit độ cũ đầu ra.
\item \path{af_detector_selection_20261031/}: protocol, selection lock, decision, dự đoán, histogram, bootstrap và audit độc lập.
\item \path{af_ecgfounder_singlelead_20261028/}, \path{af_convnext_singlelead_20261028_int8_final/}: thử nghiệm AF waveform.
\item \path{device_bench_20261023_sqi_af/}: timing/footprint SQI và AF INT8.
\end{itemize}

\section{Khả năng tái lập và kiểm tra}
Với detector-selection, \path{REPRODUCE.md} trong thư mục lượt chạy ghi thứ tự khóa protocol, kiểm tra dữ liệu, chạy sàng lọc, khóa ứng viên, chạy confirmation, audit và sinh báo cáo. Source chính gồm \path{run_detector_selection.py}, \path{run_detector_selection_shard.py} và \path{audit_detector_selection.py} trong thư mục AF experiments.

Adapter phân mảnh được dùng vì môi trường Windows chặn named pipe của ProcessPoolExecutor; các worker tính toán và quy tắc khoa học không bị thay đổi. Những lượt chưa hoàn tất hoặc lỗi được giữ riêng và không dùng thay một cohort đầy đủ.

File LaTeX này chứa toàn bộ nội dung, sơ đồ TikZ và thư mục tài liệu tham khảo, không phụ thuộc file .bib hoặc ảnh ngoài. Có thể biên dịch bằng XeLaTeX hai lần hoặc bằng Tectonic; bản cập nhật của các số liệu vẫn cần đối chiếu artifact trước khi thay trong văn bản.

\begin{thebibliography}{99}
\raggedright
\addcontentsline{toc}{chapter}{Tài liệu tham khảo}
\bibitem{pan1985} J. Pan và W. J. Tompkins, ``A real-time QRS detection algorithm,'' \emph{IEEE Transactions on Biomedical Engineering}, 32(3), tr. 230--236, 1985. \url{https://doi.org/10.1109/TBME.1985.325532}.
\bibitem{elgendi2013} M. Elgendi, ``Fast QRS Detection with an Optimized Knowledge-Based Method: Evaluation on 11 Standard ECG Databases,'' \emph{PLoS ONE}, 8(9), e73557, 2013. \url{https://doi.org/10.1371/journal.pone.0073557}.
\bibitem{laitala2020} J. Laitala và cộng sự, ``Robust ECG R-peak Detection Using LSTM,'' \emph{Proceedings of the 35th ACM/SIGAPP Symposium on Applied Computing}, 2020. \url{https://doi.org/10.1145/3341105.3373945}.
\bibitem{yun2022} D. Yun, H.-C. Lee, C.-W. Jung và cộng sự, ``Robust R-peak detection in an electrocardiogram with stationary wavelet transformation and separable convolution,'' \emph{Scientific Reports}, 12, 19638, 2022. \url{https://doi.org/10.1038/s41598-022-19495-9}.
\bibitem{orphanidou2015} C. Orphanidou, T. Bonnici, P. Charlton, D. Clifton, D. Vallance và L. Tarassenko, ``Signal-quality indices for the electrocardiogram and photoplethysmogram: derivation and applications to wireless monitoring,'' \emph{IEEE JBHI}, 19(3), tr. 832--838, 2015. \url{https://doi.org/10.1109/JBHI.2014.2338351}.
\bibitem{clifford2012} G. D. Clifford, J. Behar, Q. Li và I. Rezek, ``Signal quality indices and data fusion for determining clinical acceptability of electrocardiograms,'' \emph{Physiological Measurement}, 33(9), tr. 1419--1433, 2012. \url{https://doi.org/10.1088/0967-3334/33/9/1419}.
\bibitem{zhao2018} Z. Zhao và Y. Zhang, ``SQI Quality Evaluation Mechanism of Single-Lead ECG Signal Based on Simple Heuristic Fusion and Fuzzy Comprehensive Evaluation,'' \emph{Frontiers in Physiology}, 9, 727, 2018. \url{https://doi.org/10.3389/fphys.2018.00727}.
\bibitem{ho2025} S. Y. S. Ho và cộng sự, ``Accurate RR-interval extraction from single-lead, telehealth electrocardiogram signals,'' \emph{medRxiv}, bản thảo preprint, 2025. \url{https://doi.org/10.1101/2025.03.10.25323655}.
\bibitem{smital2020} L. Smital và cộng sự, ``Real-time quality assessment of long-term ECG signals recorded by wearables in free-living conditions,'' \emph{IEEE Transactions on Biomedical Engineering}, 2020. \url{https://doi.org/10.1109/TBME.2020.2969719}.
\bibitem{daluwatte2016} C. Daluwatte và cộng sự, ``Assessing ECG signal quality indices to discriminate ECGs with artefacts from pathologically different arrhythmic ECGs,'' \emph{Physiological Measurement}, 37(8), tr. 1370--1382, 2016. \url{https://doi.org/10.1088/0967-3334/37/8/1370}.
\bibitem{geurts2006} P. Geurts, D. Ernst và L. Wehenkel, ``Extremely randomized trees,'' \emph{Machine Learning}, 63, tr. 3--42, 2006. \url{https://doi.org/10.1007/s10994-006-6226-1}.
\bibitem{ecgfounder2025} J. Li và cộng sự, ``An Electrocardiogram Foundation Model Built on over 10 Million Recordings,'' \emph{NEJM AI}, 2(7), AIoa2401033, 2025. \url{https://doi.org/10.1056/AIoa2401033}.
\bibitem{kraft2025} D. Kraft và P. Rumm, ``Atrial Fibrillation and Atrial Flutter Detection Using Deep Learning,'' \emph{Sensors}, 25(13), 4109, 2025. \url{https://doi.org/10.3390/s25134109}.
\bibitem{shap2017} S. M. Lundberg và S.-I. Lee, ``A Unified Approach to Interpreting Model Predictions,'' \emph{Advances in Neural Information Processing Systems}, 30, 2017. \url{https://proceedings.neurips.cc/paper_files/paper/2017/hash/8a20a8621978632d76c43dfd28b67767-Abstract.html}.
\bibitem{treeshap2020} S. M. Lundberg, G. Erion, H. Chen và cộng sự, ``From local explanations to global understanding with explainable AI for trees,'' \emph{Nature Machine Intelligence}, 2, tr. 56--67, 2020. \url{https://doi.org/10.1038/s42256-019-0138-9}.
\bibitem{mecge2024} K.-H. Hung, K.-C. Wang, K.-C. Liu, W.-L. Chen, X. Lu, Y. Tsao và C.-W. Lin, ``MECG-E: Mamba-based ECG Enhancer for Baseline Wander Removal,'' \emph{IEEE International Conference on Big Data}, 2024. \url{https://doi.org/10.1109/BigData62323.2024.10825253}.
\bibitem{xie2021} Xie và cộng sự, ``A signal quality assessment-based ECG waveform delineation method used for wearable monitoring systems,'' \emph{Medical \& Biological Engineering \& Computing}, 2021. \url{https://doi.org/10.1007/s11517-021-02425-8}.
\bibitem{martinez2004} J. P. Martínez và cộng sự, ``A wavelet-based ECG delineator: evaluation on standard databases,'' \emph{IEEE Transactions on Biomedical Engineering}, 2004. \url{https://doi.org/10.1109/TBME.2003.821031}.
\bibitem{ecgdeli2021} N. Pilia, C. Nagel, G. Lenis, S. Becker, O. Dössel và A. Loewe, ``ECGdeli---An open source ECG delineation toolbox for MATLAB,'' \emph{SoftwareX}, 13, 100639, 2021. \url{https://doi.org/10.1016/j.softx.2020.100639}.
\bibitem{jaeger2024} K. M. Jaeger và cộng sự, ``Systematic Comparison of ECG Delineation Algorithm Performance on Smartwatch Data,'' \emph{IEEE Access}, 12, tr. 160794--160804, 2024. \url{https://doi.org/10.1109/ACCESS.2024.3487013}.
\bibitem{mitdb} G. Moody và R. Mark, \emph{MIT-BIH Arrhythmia Database}, phiên bản 1.0.0, PhysioNet. \url{https://doi.org/10.13026/C2F305}.
\bibitem{qtdb} R. Mark và G. Moody, \emph{QT Database}, phiên bản 1.0.0, PhysioNet. \url{https://doi.org/10.13026/C24K53}. Xem thêm P. Laguna và cộng sự, CinC 1997, \url{https://doi.org/10.1109/CIC.1997.648140}.
\bibitem{nstdb} G. Moody và R. Mark, \emph{MIT-BIH Noise Stress Test Database}, phiên bản 1.0.0, PhysioNet. \url{https://doi.org/10.13026/C2HS3T}.
\bibitem{afdb} G. Moody và R. Mark, \emph{MIT-BIH Atrial Fibrillation Database}, phiên bản 1.0.0, PhysioNet. \url{https://doi.org/10.13026/C2MW2D}.
\bibitem{ltafdb} G. Moody và S. Swiryn, \emph{Long-Term AF Database}, phiên bản 1.0.0, PhysioNet. \url{https://doi.org/10.13026/C2QG6Q}.
\bibitem{shdbaf} K. Tsutsui, S. Biton Brimer và J. Behar, \emph{SHDB-AF}, phiên bản 1.0.1, PhysioNet, 2025. \url{https://doi.org/10.13026/n6yq-fq90}. Bài mô tả dữ liệu: \url{https://doi.org/10.1038/s41597-025-04777-4}.
\bibitem{butqdb} A. Nemcova và cộng sự, \emph{Brno University of Technology ECG Quality Database}, phiên bản 1.0.0, PhysioNet, 2020. \url{https://doi.org/10.13026/kah4-0w24}.
\bibitem{cinc2011} I. Silva, G. Moody và L. A. Celi, \emph{PhysioNet/Computing in Cardiology Challenge 2011}, phiên bản 1.0.0. \url{https://physionet.org/content/challenge-2011/1.0.0/}.
\end{thebibliography}
\end{document}
